% ==============================================================
% Kapitel 07: Fazit und Ausblick
% ==============================================================
\section{Fazit und Ausblick}
\subsection{Zusammenfassung}
In dieser Diplomarbeit wurde ein Timing-System für die \gls{ZEC} entwickelt. Das System bietet eine REST-\gls{API}-Schnittstelle, die es
ermöglicht, Versuche in die Datenbank einzutragen, welche auch automatisch bewertet werden. Zusätzlich können die Ergebnisse über die \gls{API} 
abgerufen werden und alle wichtigen Informationen können einfach über die Weboberfläche eingesehen und bearbeitet werden. 

Der Microservice-Ansatz hat sich als passend für das Projekt erwiesen, da er eine klare Trennung der Verantwortlichkeiten ermöglicht und die Wartbarkeit des Systems verbessert.
Die Verwendung von Docker-Containern war ebenfalls eine richtige Entscheidung, da so die Bereitstellung erleichtert wurde und eine Unabhängigkeit von der 
Infrastruktur gewährleistet werden konnte. Die Verwendung einer zentralen Datenbank hätte die Entwicklung erleichtert, allerdings hätte dies die 
Eigenständigkeit der Services massiv eingeschränkt, weshalb die Entscheidung für einzelne Datenbanken pro Service auch hier die richtige war.

\subsection{Zukunftsausblick}

Zusätzlich zu den offensichtlichen Erweiterungen, wie der Fertigstellung der noch nicht implementierten Features, etwa der Desktop-App zur Verbindung mit 
der Lichtschranke per MQTT, der Kubernetes-Integration und den \gls{CICD}-Pipelines, kamen während der Entwicklung auch andere Ideen auf, die in Zukunft umgesetzt werden könnten.

\begin{itemize}
    \item \textbf{SD-Karten-Schutz: } Derzeit ist kein Schutz für die SD-Karten der Raspberry Pis implementiert. Eine Implementierung von Caching im Backend oder Frontend
    wäre notwendig, sollte der Server auf einem Raspberry Pi laufen.
    \item \textbf{Anbindung einer Bluetooth-Datenübertragung: } Im derzeitigen System muss der Zeitnehmer von den Sensoren ablesen, wie viel Energie verbraucht wurde.
    Eine Anbindung der Bluetooth-Datenübertragung könnte es ermöglichen, dass die Sensoren die Daten an die Desktop-App übermitteln.
    
\end{itemize}
